\chapter*{General conclusion}
\addcontentsline{toc}{chapter}{General conclusion}
\markboth{General conclusion}{General conclusion}
\label{chap.general-conclusion}

Inventory is among the most repetitive and least automated operations in
warehouse management. Aerial robots are particularly well suited to this task, as
they can reach every storage level and read the labels already attached to the
goods. Deploying a team of drones, however, raises challenges that extend well
beyond label reading: the team must achieve full coverage of the building,
allocate the work among its members, and remain reliable when the mission
deviates from its expected course.

In the first part of this report, we introduced the foundations required to
understand this problem. We first presented warehouse inventory, its operational
cost, and the use of aerial robots in logistics. We then reviewed existing
approaches to autonomous inventory, multi-robot exploration and task
coordination. This review showed that most drone-based inventory systems depend
on a map of the warehouse known in advance, from which vehicle routes and task
allocation are fixed before take-off. Such systems perform efficiently as long as
the mission follows the plan, but they offer no mechanism to cope with a vehicle
failure, a change in the environment during the flight, or the absence of a prior
map.

In the second part, we presented SwarmScan, a system in which a team of three
drones performs the inventory of a warehouse without any prior knowledge of its
layout. The vehicles incrementally build a shared map of the building during the
flight, and each of them selects its next target from the collective knowledge of
the team, without relying on a central planner. The system was implemented in a
simulation environment that generates warehouses procedurally and controls each
vehicle through a realistic autopilot, thereby ensuring that every result stems
from a complete flight.

The implemented system introduced several mechanisms to ensure the completeness
and robustness of the inventory. The shared map aggregates the explored space,
the identifiers already read and the labels detected but not yet read, so that
every decision relies on the knowledge of the entire team. A claim mechanism
allows each vehicle to signal the target it is pursuing, thereby preventing
duplicated effort, while claims expire automatically when a vehicle ceases to
operate. A learned detector identifies labels at a distance, allowing the team to
direct its effort before the labels come within reading range. In addition, a
semantic guide based on a \gls{abb:vlm} provides high-level advice on the areas
to explore next.

The experimental evaluation demonstrated the effectiveness of this approach.
SwarmScan read 96.5\% of the identifiers in the reference warehouse and 92.3\% in
a warehouse never encountered during its design, without any adjustment. Compared
with a static approach that plans the entire mission from a known map, it reached
90\% of the inventory more than two and a half times earlier, and required 1.9 to
2.5 times less flight distance to achieve the same result. The evaluation also
highlighted the robustness of the system. The loss of a vehicle during the
mission had no effect on the final inventory, whereas the static approach lost 31
identifiers under the same conditions. The sudden appearance of an obstacle was
handled without the loss of a single identifier. Across all missions, no vehicle
entered a rack, no target was pursued by two vehicles simultaneously, and no
non-existent identifier was ever recorded.

These results provide answers to the four research questions of this work. First,
a team of drones can decide where to go and what to read by relying solely on the
knowledge acquired during the flight, as the map it builds is accurate enough to
serve as the sole basis for its decisions. Second, the work can be distributed
without a central planner, since the claim mechanism proved sufficient to prevent
duplicated effort and collisions in every mission. Third, the team remains robust
to disruptions, as the work of a lost vehicle is naturally redistributed to the
remaining ones and obstacles are integrated into the map as soon as they are
perceived. Finally, a team that starts without any map outperforms a static
approach with full prior knowledge of the building, and the margin widens
precisely in the situations that the plan fails to anticipate.

Overall, this work demonstrates that prior knowledge of the warehouse layout is
not a prerequisite for reliable drone-based inventory. A team that builds its own
representation of the environment and continuously adapts its decisions to it
proves both more efficient and more robust than one that executes a plan fixed
before take-off.

\section*{Discussion and Future Work}

Although SwarmScan demonstrates the feasibility of map-free inventory by a team
of drones, several directions remain open. The current system has been developed
and evaluated exclusively in simulation, and it does not read the smallest
labels, which would require operating closer to the shelves than the system
allows. This work therefore provides a first complete foundation, whose
deployment in real warehouses calls for several extensions.

\textbf{Validation on Physical Drones}: The autopilot used in simulation is the
same software that runs on real vehicles, which facilitates the transfer of the
system to physical platforms. An experimental campaign in a real warehouse would
validate the detection and reading performance under real lighting conditions,
motion blur and sensor noise.

\textbf{Reading the Smallest Labels}: The identifiers missed by the system are
carried by labels approximately three times smaller than the others. A dedicated
reading strategy, triggered for labels that are detected but not decoded, could
approach them at a shorter distance and a lower speed in order to complete the
inventory.

\textbf{Dynamic Warehouses}: The evaluation showed that the system integrates an
obstacle appearing during the mission. Operational warehouses, however, are
highly dynamic environments, in which forklifts and operators circulate in the
aisles and the stock evolves throughout the day. Extending the shared map to
moving obstacles and to changes in the stored goods would bring the system closer
to deployment in active facilities.

\textbf{Scalability to Larger Teams and Buildings}: Since no vehicle plays a
central role, increasing the number of drones does not alter the decision
process. Evaluating the system with larger teams, in larger warehouses and across
a wide range of generated layouts, would assess the limits of this scalability
and provide statistically grounded results.

Overall, these perspectives position the current work as a first operational step
toward autonomous inventory by teams of drones in real warehouses. The system
establishes the core capabilities required for this objective, namely building a
map during the flight, making decisions without a central planner, and remaining
robust to unexpected events. Future work will build on this foundation to
progressively bring the system from simulation to real-world deployment.
